\documentclass[12pt, a4paper]{article}

\usepackage[utf8]{inputenc}
\usepackage[T1]{fontenc}
\usepackage{lmodern}
\usepackage[margin=1.2in]{geometry}
\usepackage{graphicx}
\usepackage{hyperref}
\usepackage{microtype}
\usepackage{booktabs}
\usepackage{array}
\usepackage{caption}

\hypersetup{
  colorlinks=true,
  linkcolor=black,
  citecolor=black,
  urlcolor=black
}

\title{Wynn's World: The Cognitive Architecture at Work\\
in a Minimal Counting Scenario}

\date{}

\begin{document}

\maketitle

\begin{abstract}
The architecture proposed in \cite{architecture} (hereafter `the architecture
paper') is implemented and tested in a minimal scenario inspired by Wynn's
(1992) infant counting experiments. While both the mechanism and the
interpretations of Wynn's findings remain debated, the experimental setup ---
controlled perception, known arithmetic, measurable violation --- is a tight
fit for a focused implementation test of the proposed architecture.

A brute-force matcher searches for productive wirings between an unsupervised
perception stream and pre-trained GRU-based predictive models, using a deficit
scoring function as the sole evaluation criterion. The implementation produces
results that were not designed in: underdetermination of wirings is
encountered, not stipulated; the proto-ontology is required, not implemented
as a feature; model succession emerges as a rank crossing in deficit tables,
not as a scripted transition. We report multiple experimentally confirmed
findings, including the distinction between trained and untrained spare
capacity, two qualitatively different environmental signals, lateral detection
of model inadequacy before new tokens arrive, succession as an inspectable
architectural event, and robustness of all findings across scoring functions.
All experiments are reproducible from the accompanying repository \cite{repo}.
\end{abstract}

\section{Introduction}
\label{sec:intro}

The architecture paper proposes a cognitive architecture built on the
independence of perception and prediction, with a matching process as the sole
bridge between them. This paper puts that architecture to work.

The immediate question is: is this architecture demonstrable at all? Can one
define an experimental setup basic enough to be tractable, but not so basic
that the results are trivial?

This is a difficult call. Luckily, we can follow the footsteps of some
brilliant behavioral researchers.

In 1992 Karen Wynn was working as a post-doc at the University of Arizona. She
designed an experiment methodologically ingenious both in the arrangement and
the measurement.

In a dark room a 5-month old infant would be sitting, facing a single bright
object: the stage. Behind them, out of sight, their parents would be present,
providing the feeling of security, while not contaminating the infant's
perception with social interactions. On stage, equipped with a rotating screen,
Mickey Mouse dolls would appear and disappear, going behind and coming out from
behind the screen. Sometimes the screen is moved away, revealing how many dolls
are at this moment behind it.

First the experiment goes through a familiarization phase. The world does not
cheat: if two dolls have gone behind the screen and the screen goes up, two
dolls are revealed. If two dolls were behind the screen, and one of them has
been seen going away, removing the screen reveals one doll, as expected.

Then the world begins to --- occasionally --- cheat. The count of the dolls
behind the screen defies the rules of arithmetic.

The measurement is amazingly simple: it is the temporal trace of infant's
attention, or, in plain words, how long it looks at the revealed dolls. The
correlation between the attention given and the arithmetical mismatch is the
experiment result: five-month-olds, who can't speak, can't point, can barely
move, demonstrate through looking time alone that something in their head
expected a different count.

We are going to adopt this setup almost unchanged: the unmodulated perception
pipeline, the perception signals relating only to dolls arrivals, departures,
and their observed count. The only adaptations we are making are increasing the
number of dolls to better register the interesting phenomena within the
architecture.

This is a companion paper. Its goal is not to demonstrate cognitive realism.
The goal is to show that the architecture presented in \cite{architecture},
taken seriously as an implementation constraint, forces outcomes that the
architecture paper claims. Where the architecture paper argues, this paper
exhibits.

\section{The Minimal World}
\label{sec:world}

\textbf{World dynamics.} Perception needs some external observables to work
on. Following Wynn's elimination and focus, the ``world'' in our experiment has
two events and one state, corresponding to doll arrival, doll departure, and
dolls count, respectively.

More precisely: the world maintains an integer count state. At each step, one
of three things happens: an increment (INC), a decrement (DEC), or a read of
the current count (C0, C1, C2, \ldots). The world deals with a limited number
of dolls available. If there are three dolls available, the possible count
readouts are (C0, C1, C2, C3). We call such a world a ``capacity 4 world'' or,
in short, cap4 world.

The world is realistic: the observed number of dolls never turns negative, nor
does it surpass the total number of dolls available. DEC cannot happen if the
current count is C0; INC is not possible if the count is at the maximum.

Otherwise the world works as expected, emitting an INC event when increasing
the state, and a DEC upon decreasing it.

Following the correspondence to the real world setup, the ``capacity'' of the
world is a parameter that can be changed: another doll can be introduced,
moving this limit of the world and making another readout available.

A cap3 world has count states \{0, 1, 2\} and token vocabulary \{INC, DEC, C0,
C1, C2\}. A cap4 world has count states \{0, 1, 2, 3\} and token vocabulary
\{INC, DEC, C0, C1, C2, C3\}. A cap5 world extends further.

The infant in the dark room observes the scene. If it observes a doll going
in --- we model this as an INC perception token. A doll goes out --- a DEC.
Observing dolls on the scene results in different tokens C0, C1, \ldots,
depending on the dolls count. These are not ``counts'' --- they are just
distinct tokens. We label them the way we do because \textit{we} know. The
infant does not. At least not initially.

There is nothing more in this world. The infant is concerned only with tracking
numerosity of a single category.

\section{Implementation}
\label{sec:implementation}

\textbf{Boxed predictive models.} The architecture calls for perception to be
contrasted with and played against a portfolio of predictive automata. In our
implementation each automaton is a GRU-based recurrent network, pre-trained in
isolation on sequences from its training world. The model exposes a fixed-size
token dictionary and returns probability distributions over next tokens. It is
important that the dictionary and the probability distribution are exposed
uniformly: even if internally some tokens correspond to state transitions, and
others to state readouts, this distinction is not published in any way outside
the model. The matcher sees only the dictionary size and the output
distributions.

The full encapsulation of the internal predictive logic of a model allows for
inclusion of a fully deterministic symbolic counter as an interesting benchmark
and a useful sanity check.

\textbf{Curriculum training.} Models are trained in stages:
cap1$\to$cap2$\to$cap3$\to$cap4, each stage consolidating before the next.
Without curriculum, the GRU achieves similar accuracy but with internally
scrambled count state representations.

Three model variants are used: ``native,'' with full curriculum pre-training;
``ghost,'' with pre-training stopped early; and ``symbolic'' --- a
deterministic benchmark.

Most of the actual experiments use three models:

\begin{center}
\begin{tabular}{llll}
\toprule
Model & Architecture & Training & Purpose \\
\midrule
acc70 & cap4 GRU & Full curriculum, $\sim$70\% accuracy & Lower-quality baseline \\
acc80 & cap4 GRU & Full curriculum, $\sim$80\% accuracy & Primary experimental model \\
ghost & cap4 GRU (or cap5) & Curriculum stopped early & Tests adaptive potential \\
\bottomrule
\end{tabular}
\end{center}

\textbf{The matcher.} Following the architecture paper, the system must find
the wirings between the perception stream and the predictive models that result
in low predictive error. In this experiment this is done brute-force /
incremental: initially the matcher evaluates every possible wiring (permutation
of perception tokens to model tokens) against the perception stream. Each
wiring is scored by deficit: the cumulative shortfall between the model's
predicted probability for the actual next token and the uniform baseline. Upon
expansions of the working dictionaries, the matching process evaluates possible
extensions of its current portfolio of wirings.

The matcher's internal state --- the structure that makes it possible to ask
``what if perception token X maps to model token Y?'' before any mapping
exists --- is what the architecture paper calls the proto-ontology. In this
implementation, it is the mapping vector between perception tokens and model
tokens: initially unassigned, populated by the matching process, and retained
as the fossilized wiring. The matcher is stateful not only for the sake of the
matching process: retaining the predictively performant wirings is necessary
both for exploitation and for the incremental adaptation.

The brute-force search is feasible here because the world is small. It is not a
requirement of the architecture --- it is a choice that makes the matcher's
behavior fully inspectable.

\textbf{Deficit and the renormalized distribution.} The predictive deficits of
the available wirings must be calculated in order to identify the best ones.
The architecture paper calls for a strict isolation between the internal
machineries of the perception, the prediction, and the matcher. This makes the
design of the deficit function an interesting exercise. The function must be
oblivious to the internal ontologies of the models, should be independent of
history, and should be simple.

Using the uniform distribution over the model's token dictionary as a baseline,
we define the predictive error at any point as $\max(1/n_{\mathrm{tokens}} -
p_{\mathrm{pred}}(\mathrm{actual\_next}),\ 0)$, and the predictive deficit as
the error summed over steps, averaged over sequences. A wiring that predicts no
better than chance accumulates maximum deficit; a wiring whose model
confidently predicts the correct next token accumulates zero.

The choice of this specific deficit formula affects what is visible but not
what is true. As shown in Appendix~\ref{app:findings}, the architectural
findings hold across scorer families, though different scorers expose different
aspects of the deficit landscape.

\textbf{Cliff pruning.} After scoring all wirings, contenders are selected by
finding the largest gap in the sorted deficit list. Everything below the gap is
retained. This is adaptive --- it reads the structure of the data rather than
imposing a fixed threshold. The contender count is the matcher's epistemic
humility about which wiring will generalize.

\textbf{Incremental wiring.} When the world expands (e.g., cap3$\to$cap4),
the matcher does not restart from scratch. It extends each surviving contender
by searching only over the unassigned model tokens for the new perception
token. Fossilized wirings constrain the search space.

\textbf{Experimental conditions.} Unless otherwise noted: n\_seeds=5,
seq\_len=300, seed=42, p\_read=0.6.

\section{The Matcher Finds Correct Wirings}
\label{sec:matcher}

We begin the experiment with the world that contains two dolls. At most two
dolls are visible at any moment. Three dolls visible are not possible --- there
are not that many dolls in this world. The world does not cheat: the arithmetic
is truthful, the number of dolls does not decrease below zero nor does it
increase over two.

The infant tries to align its predictive potential with perception. In Wynn's
terms this corresponds to \textit{familiarization}.

The matching process arrives at the following ranking of the possible wirings:

\begin{center}
\begin{tabular}{rlrrrrrr}
\toprule
Place & Deficit & Model & INC & DEC & C0 & C1 & C2 \\
\midrule
1  & 0.0000 & cap4\_ghost  & INC & DEC & C0 & C1 & C2 \\
2  & 0.0102 & cap3\_acc70  & INC & DEC & C0 & C1 & C2 \\
3  & 0.0136 & cap4\_ghost  & DEC & INC & C2 & C1 & C0 \\
4  & 0.0287 & cap3\_acc80  & DEC & INC & C2 & C1 & C0 \\
5  & 0.0573 & cap4\_acc80  & DEC & INC & C2 & C1 & C0 \\
6  & 0.0712 & cap3\_acc70  & DEC & INC & C2 & C1 & C0 \\
7  & 0.0789 & cap3\_acc80  & INC & DEC & C0 & C1 & C2 \\
8  & 0.2405 & cap4\_acc80  & INC & DEC & C0 & C1 & C2 \\
9  & 0.3025 & cap4\_acc80  & INC & DEC & C0 & C1 & C3 \\
10 & 0.3857 & cap4\_acc80  & DEC & INC & C3 & C1 & C0 \\
11 & 1.5831 & cap3\_acc70  & C1  & DEC & C0 & INC & C2 \\
12 & 1.6048 & cap4\_acc80  & INC & DEC & C1 & C2 & C3 \\
\bottomrule
\end{tabular}
\end{center}

There are three things to observe about this table.

\textbf{The Cliffs.} We have truncated it at the twelfth position to make the
cliffs visible. The first cliff is between position 7 and position 8. The
second cliff, between positions 10 and 11, is an order of magnitude steeper.
Any wiring up to the first cliff is a reasonable contender. Any wiring beyond
the second cliff can be discarded. The fate of the wirings in-between depends
on the system's appetite for adaptability.

\textbf{The Symmetry.} There are two families of wirings which perform well:
the ``direct'' wirings, and ``mirror-image'' wirings. The direct wirings are
what we expected: the world's INC event is mapped to the model's INC, DEC to
DEC, and the numerosities are mapped directly. The mirror-image wirings are
only obvious in hindsight: the world has a reflection symmetry which the mapper
faithfully reports. After truncating the contenders list at the cliff, the
system is holding two contradictory-but-equal bets about where the world's
count axis sits in the model's vocabulary.

\begin{center}
\begin{tabular}{lrrrrrr}
\toprule
Wiring family & INC$\to$ & DEC$\to$ & C0$\to$ & C1$\to$ & C2$\to$ & Best deficit \\
\midrule
Canonical         & INC & DEC & C0 & C1 & C2 & 0.0000 \\
Reverse-canonical & DEC & INC & C2 & C1 & C0 & 0.0136 \\
\bottomrule
\end{tabular}
\end{center}

In the context of architecture paper claims: the matcher does not care for nor
is aware of the internal logic of a model. Operational performance is all the
signal. The architectural principle has determined what has unfolded.

\textbf{The Ghost.} The best performing wiring connects to the model with an
\textit{untrained} C3 pad. The perfect performance of this model in the cap2
world is a surprise, and can probably be attributed to its internal spare
capacity. As demonstrated in the next section, however, betting solely on this
model in a world that may evolve would be perilous.

\textbf{Note on the deficit function.} The specific formula we have chosen
involves a very simple renormalization of the returned probability distribution,
which depends on the wired dictionary size. Without renormalization, a model
that spreads mass across unwired pads would appear worse than it is. This is
the only place in the architecture where any calibration of a ``capN model''
to a ``capM world'' happens.

\section{The Arrival of the Third Doll}
\label{sec:thirddoll}

\textbf{Incremental familiarization.} It had taken approximately five months
for the infants from Wynn's experiment to arrive at a developmental stage where
the mismatch on counts up to 2 would trigger the observed behavior. We have no
idea how much more time it would take them to display similar behavior with up
to 3 dolls. For us, however, this is a very natural next step in the
experiment.

When the world expands incrementally (cap3$\to$cap4$\to$cap5), fossilized
wirings from the previous stage constrain the search at the next stage. The
effect is dramatic.

\begin{center}
\begin{tabular}{llll}
\toprule
Stage & World & Evaluations per model & Brute-force equivalent \\
\midrule
Stage 1 & cap3 (5 tokens) & $P(7,5) = 2{,}520$ & --- \\
Stage 2 & cap4 (6 tokens) & $9 \times 2 = 18$ contenders & $P(7,6) = 5{,}040$ \\
Stage 3 & cap5 (7 tokens) & $14 \times 1 = 14$ contenders & $P(7,7) = 5{,}040$ \\
\bottomrule
\end{tabular}
\end{center}

At stage 2, instead of evaluating 5,040 wirings from scratch, the matcher
extends 9 surviving contenders over 2 remaining model tokens: 18 evaluations.
At stage 3, 14 extensions with 1 remaining token each. The fossilized wirings
are not just records --- they are search constraints. This is the computational
consequence of fossilization that the architecture paper describes as ``the
system's ontological commitments.''

The incremental strategy works only as long as the final contenders survive the
pruning at the earlier stages. We found that our simple ``cliff'' heuristic
performs well. We do not go into detailed study of the cutoff regimes to keep
this paper focused.

\textbf{Adaptive potential.} We model the arrival of the third doll --- the
evolution of the world --- as a simple change to the perception token
sequences: we include C3 as a new token in the perception vocabulary, and we
move the constraint on the appearance of INC from C2 to C3 (allowing for
C2$\to$C3 transitions and forbidding increases beyond C3). We run the matching
competition on those contenders from the cap3 pool which have what the
architecture paper calls adaptive potential (i.e., at least one unwired pad).

In this pool we have a ``native'' model, pre-trained to count up to three, and
a ``ghost'' model. The ghost model has the same architecture as a native model
(same number of tokens, same GRU) but its curriculum was stopped early. A cap4
ghost trained only to cap3 has the C3 slot in its vocabulary, but C3 received
zero gradient during training. The slot exists architecturally but carries no
representation. Formally, until tested against the cap4 world, both models
promise the same adaptive potential.

\textbf{On the home world (cap3), the ghost wins:}

\begin{center}
\begin{tabular}{rlll}
\toprule
Place & Deficit & Model & Wiring \\
\midrule
1 & \textbf{0.0000} & cap4\_ghost  & canonical \\
2 & 0.0136          & cap4\_ghost  & reverse-canonical \\
3 & 0.0577          & cap4\_acc80  & reverse-canonical \\
4 & 0.3795          & cap4\_acc80  & canonical \\
\bottomrule
\end{tabular}
\end{center}

The ghost's relaxed interior representation (cap4 architecture trained to cap3
boundary) outperforms native cap3 models on their own world.

\textbf{On the extended world (cap4), the ghost collapses:}

\begin{center}
\begin{tabular}{rllll}
\toprule
Place & Deficit & S1 Base & Model & Wiring \\
\midrule
1 & \textbf{0.0400} & 0.3795 & cap4\_acc80  & canonical \\
2 & 3.1905          & 0.5072 & cap4\_acc80  & shifted \\
3 & 7.0694          & 0.0000 & cap4\_ghost  & canonical \\
4 & 8.6570          & 0.0136 & cap4\_ghost  & reverse-canonical \\
\bottomrule
\end{tabular}
\end{center}

The ghost's leading wiring deficit goes from 0 to over 7. The ghost's C3 slot
is inert. No valid wiring exists --- all ghost wiring attempts on cap4 score
$\geq 6.36$.

We see a dramatic example of what the architecture paper calls wiring and model
succession. The value of retaining contenders is demonstrated: the correct path
forward was not the best current wiring but a weaker contender with the right
base for extension.

\textbf{Wiring succession (to the same model).} The architecture paper demands
that we demonstrate succession between competing wirings to the same model.
This is clearly visible in the scores of the cap4\_acc80 model against cap3
and cap4 worlds:

\medskip
\noindent\textit{2-Dolls World}
\begin{center}
\begin{tabular}{rlrrrrr}
\toprule
Deficit & Model & INC & DEC & C0 & C1 & C2 \\
\midrule
0.0573 & cap4\_acc80 & DEC & INC & C2 & C1 & C0 \\
0.2405 & cap4\_acc80 & INC & DEC & C0 & C1 & C2 \\
\bottomrule
\end{tabular}
\end{center}

\noindent\textit{3-Dolls World}
\begin{center}
\begin{tabular}{rlrrrrrr}
\toprule
Deficit & Model & INC & DEC & C0 & C1 & C2 & C3 \\
\midrule
0.0400 & cap4\_acc80 & INC & DEC & C0 & C1 & C2 & C3 \\
0.9823 & cap4\_acc80 & DEC & INC & C3 & C2 & C1 & C0 \\
\bottomrule
\end{tabular}
\end{center}

The mirror-image wiring performs better in the smaller world, but collapses
when the world extends. The model remains, but a different wiring takes the
lead.

\section{Two Kinds of Environmental Change}
\label{sec:twokinds}

Intuitively, there is a difference for an infant between observing something
\textit{new} and observing something \textit{wrong}. The very reason for the
familiarization phase in Wynn's experiment is to remove the \textit{new} and
isolate the \textit{wrong}. In our simple setup, this intuition becomes
tractable. The architecture paper describes three signals ruling the system's
mechanics: the predictive error, the adaptive potential, and the presence of
unwired perception tokens (which the paper calls the ``cognitive itch'').

Let us narrow our focus to the two models from the previous sections: the
``native'' cap4\_acc80, and the cap4\_ghost. The cap3 world has been
``familiarized'': optimal wirings to both models have been identified and
fossilized (see the score table in Section~\ref{sec:matcher}).

In the previous section we have seen the outcome of the re-wiring effort. Now
let us inspect in some detail the ontological crisis as it unfolds, but before
it is overcome. What are the signals available when the world changes, which
could guide the adaptation?

\textbf{The New.} The infant has familiarized the world of two dolls. The
parent, or the nice researcher, has brought another doll. The world, as
reported by the perception stream, offers new phenomena: there is now a C3
perception token available, and INC can happen when the count state is 2. How
do the wirings fossilized at cap3 perform?

It turns out that even if the new token is filtered from the evaluation stream,
the ghost still fails. We verified this directly: filtering all C3 tokens from
the evaluation stream, the ghost still fails with a deficit of 3.3845 --- nearly
half of its unfiltered deficit of 7.0694. The boundary dynamics alone, without
any new token present, are sufficient to expose model inadequacy. Where the
architecture paper talks about ``blindness'' to the arrival of new tokens, we
identify a ``lateral detection'' of the world change. Its specific form is
determined by the nature of our experimental setup, but in general we expect
this phenomenon to stem from the predictive nature of the models and to occur
often.

A simple simulation shows how the prediction deficits change when the world
extends, but the wirings remain:

\begin{center}
\begin{tabular}{rrrrrrrr}
\toprule
Place & Old Deficit & New Deficit & Model & INC & DEC & C0 & C1 \\
\midrule
1 & 0.3795 &  0.0400 & cap4\_acc80 & INC & DEC & C0 & C1 \\
2 & 0.0577 & 10.3050 & cap4\_acc80 & DEC & INC & C3 & C2 \\
3 & 0.0000 &  7.0694 & cap4\_ghost & INC & DEC & C0 & C1 \\
\bottomrule
\end{tabular}
\end{center}

The one wiring of the native counter performs well, even significantly
improving its performance --- it correctly predicts the INC/DEC dynamics at
the new boundary without observing the intermediate count state. The additional
world structure disambiguates this wiring. The inverse wiring collapses. The
ghost counter collapses too --- it has never seen INC events happening when the
count is C2.

Decomposing the deficit by predicted token type reveals a directed gradient:

\begin{center}
\small
\begin{tabular}{lrrrrrrr}
\toprule
Wiring & INC & DEC & C0 & C1 & C2 & C3 & Total \\
\midrule
ghost canonical  & 0.0000 & 0.0000 & 0.0000 & 0.0300 & 3.3545 & 3.6849 & 7.0694 \\
/step   & 0.0000 & 0.0000 & 0.0000 & 0.0001 & 0.0120 & 0.0297 &        \\
ghost reverse    & 0.0000 & 0.0000 & 0.0136 & 0.2802 & 4.3211 & 4.0422 & 8.6570 \\
/step   & 0.0000 & 0.0000 & 0.0001 & 0.0008 & 0.0154 & 0.0326 &        \\
acc80 canonical  & 0.0000 & 0.0000 & 0.0000 & 0.0400 & 0.0000 & 0.0000 & 0.0400 \\
/step   & 0.0000 & 0.0000 & 0.0000 & 0.0001 & 0.0000 & 0.0000 &        \\
acc80 reverse    & 0.0000 & 0.0000 & 0.0549 & 1.1345 & 5.0430 & 4.0725 & 10.3050 \\
/step   & 0.0000 & 0.0000 & 0.0004 & 0.0031 & 0.0180 & 0.0328 &        \\
\bottomrule
\end{tabular}
\end{center}

The results are telling: the error accumulates at the C2/C3 boundary, except
for the one wiring of the native model that happens to be aligned with the
extension. The other wiring of the native model performs the worst. The ghost
model fails in a more symmetric fashion. It tolerates orientation reversal
because it has nothing at C3 to lose; the native model's wirings split because
it does.

The lateral detection provides a structured fingerprint of the world change,
and is accompanied by the ``itch'' of the arrival of a new, unwired perception
token. It turns out that the infant has surprisingly rich data prior to
initiating the search for productive extensions of its fossilized wirings.

\textbf{The Wrong.} Let us now consider the situation where the world starts
``cheating.'' There is nothing new in the world or in the perception stream,
but 10\% of the world read events report the wrong count (C$_k$ replaced by a
random C$_j$, $j \neq k$). The ranking of the prediction deficit of the
wirings in the portfolio changes:

\begin{center}
\begin{tabular}{rrrlrrrrr}
\toprule
Place & S1 Def & Viol Def & Model & INC & DEC & C0 & C1 & C2 \\
\midrule
1 & 0.0136 & 1.1062 & cap4\_ghost & DEC & INC & C2 & C1 & C0 \\
2 & 0.0000 & 1.1137 & cap4\_ghost & INC & DEC & C0 & C1 & C2 \\
3 & 0.0577 & 1.3079 & cap4\_acc80 & DEC & INC & C2 & C1 & C0 \\
4 & 0.3795 & 1.6938 & cap4\_acc80 & INC & DEC & C0 & C1 & C2 \\
\bottomrule
\end{tabular}
\end{center}

In contrast with the world extension, now all wirings' predictive scores
degrade. There are no obvious candidates for a wiring extension. Moreover,
there are no new tokens in the perception stream. For our simple system the
only way to improve the predictive performance is to reset the predictive
portfolio and re-bootstrap. A system with a more elaborate perception pipeline
might try to focus on identifying any missing perception details. This is, in
our picture, the origin of the ``heightened attention'' of the infants in
Wynn's experiment --- something more basic than ``surprise.'' Under the
experiment's conditions there is no productive adaptation path available. What
remains is re-processing.

\textbf{The distinction.} The world extension and the screen violation are
therefore distinguishable not just in magnitude but in architecture: extension
produces a directed signal that concentrates at the training boundary and
discriminates in favour of the model with trained capacity. Violation produces
a uniform, inverted signal that carries no information about which model should
succeed. The signatures are very different. The baby can in principle
distinguish between seeing something new and being fooled.

Note: these patterns repeat themselves when the parent brings yet another doll
--- as we have tested for the cap4$\to$cap5 transition.

\section{Discussion}
\label{sec:discussion}

The original purpose of the paper was to provide a toy model of the cognitive
architecture described in \cite{architecture}. Modelling the streamlined world
of Wynn's experiment has allowed us to capture the architecture in a very
mechanistic form. Experimenting with actual implementations, however simple,
brought to light a few new aspects, and refined our views --- still staying
truthful to the architecture paper. The architecture paper argues that certain
outcomes must follow from certain principles. This paper puts those principles
to work and reports what happened.

The underdetermination wasn't stipulated --- it appeared in the deficit table
at stage 1, where the ghost and its mirror image both scored near zero. The
succession wasn't programmed --- it fell out of comparing two ranked deficit
lists across stages. The findings are the architecture paper's claims; the
implementation is their witness. Appendix~\ref{app:findings} contains a full
list of observations and findings that had not been pre-programmed by the
architecture paper.

\textbf{What the implementation forced.}

\textit{The proto-ontology.} The matching process requires a working state: a
data structure that records, for each perception token, which model token it is
currently mapped to. This was not designed in as a named architectural
component. It was required because the incremental wiring strategy needed
somewhere to record which assignments were fossilized. The architecture paper
argues that proto-ontology is a requirement of any perception-prediction
coupling architecture. The implementation confirms: there is no way to
implement incremental wiring without something that plays this role. Without
retaining the top wirings fossilized in the proto-ontology (i.e., the
matcher's state), the whole matching process would be pointless, and
incremental adaptation impossible.

\textit{The interface principle} is confirmed as a functional requirement: when
violated, the matching problem becomes unsolvable
(Appendix~\ref{app:findings}, Finding~13).

\textit{The ghost result.} The architecture paper distinguishes architectural
capacity from trained capacity. The ghost has a slot for C3 but no trained
representation behind it. This distinction was expected. What the
implementation reveals is the magnitude: a deficit gap of 7.03 between ghost
and native on the extended world, while the gap on the home world runs in the
other direction. The numerical gap is not a design choice --- it reflects how
severely the GRU's hidden state is corrupted when an untrained slot is forced
into active use. The architecture paper describes this as a prediction-driven
discrimination signal. The 7.03 is that signal's amplitude.

\textit{Lateral detection.} This was not predicted by the architecture paper
and was not a design goal. The implementation produced it as a side effect of
routing all signals through a single deficit scorer. The ghost's failure is
visible before the new token appears because the INC/DEC dynamics at the
training boundary are already wrong. As described in
Section~\ref{sec:twokinds}, this was discovered by running the experiment, not
by reasoning about it in advance. The architecture paper describes the
cognitive itch (Signal~2) as the primary trigger for wiring search. The
implementation adds a prior signal --- predictive deficit degradation appears
before the cognitive itch. The architecture's self-diagnosis is earlier than
the architecture paper describes.

\textbf{What remains open.}

The world is minimal. Three to five count states, GRU-based models, brute-force
search. The architecture paper claims at Marr's computational level; this paper
exhibits at the algorithmic and implementation levels for one scenario. Whether
the findings generalize to richer worlds, larger models, and approximate
matching is an empirical question.

Whether the reproducible findings in a minimal world constitute theory
confirmation or illustration is the appropriate skeptical question. The
architecture paper's own criterion is that the architecture should make
observable predictions in implemented systems. These findings are that
observation. The reader is in the better position to judge whether the
threshold has been crossed.

The brute-force matcher is the control condition. It establishes what the
architecture can find. Approximate matchers are the scaling question, but you
cannot evaluate them without first knowing what the correct answer is.

\section{Conclusion}
\label{sec:conclusion}

The implementation forces things. The proto-ontology is forced by the need for
an incremental state. The contender pool is forced by the mathematics of world
extension --- the correct model cannot always be the current leader. The
succession is forced by the same scoring function that tracked performance on
the home world. Lateral detection is forced by the GRU's boundary dynamics.
None of these were designed. They were encountered.

The architecture paper argues these outcomes follow from three independence
principles: perception doesn't know the prediction vocabulary, prediction
doesn't know the perception stream, and matching evaluates the coupling without
access to either interior. This paper implements those principles and reports
what follows. The cliff heuristic embodies the specific epistemic stance the
architecture requires: intermediate scores are unreliable discriminators, and
preserving minority contenders is the cost of remaining open to a world that
has not yet finished changing.

The code and all experiments are available in the accompanying repository
\cite{repo}.

%% ============================================================
\appendix
\renewcommand{\thesection}{Appendix \Alph{section}}
%% ============================================================

\section{Complete Findings Register}
\label{app:findings}

The following findings emerged from the implementation. Items marked [body] are
presented in the main text; items marked [appendix] are documented here with
supporting data. All are reproducible from the repository.

\subsection*{Architectural findings}

\noindent\textbf{F1. The matcher recovers correct wirings.}
Brute-force search over all permutations identifies the two structurally valid
wiring families (canonical and reverse-canonical) across all models and seeds.
[body, Section~\ref{sec:matcher}]

\medskip
\noindent\textbf{F2. Underdetermination is encountered, not stipulated.}
The two valid families reflect the counting world's reflection symmetry. The
matcher reports this faithfully --- it was not designed to find symmetry.
[body, Section~\ref{sec:matcher}]

\medskip
\noindent\textbf{F3. The deficit landscape has cliff structure.}
Wirings cluster into quality tiers separated by sharp gaps. The cliff between
valid and invalid wirings is an order of magnitude steeper than the gap between
valid families. [body, Section~\ref{sec:matcher}]

\medskip
\noindent\textbf{F4. Fossilization collapses the search space.}
Incremental wiring reduces stage-2 search from $P(7,6) = 5{,}040$ evaluations
to 18, and stage-3 from $P(7,7) = 5{,}040$ to 14. Fossilized wirings are
search constraints, not just records. [body, Section~\ref{sec:thirddoll}]

\medskip
\noindent\textbf{F5. Architectural spare capacity without trained knowledge is
not adaptive potential.}
The ghost model (same architecture, curriculum stopped early) wins on the home
world but collapses on the extended world with a deficit gap of 7.03. The slot
exists but carries no representation. Replicated at cap5 scale.
[body, Section~\ref{sec:thirddoll}]

\medskip
\noindent\textbf{F6. Model succession is an inspectable rank crossing.}
The ghost leads at stage 1 (deficit 0.0000); the native model leads at stage 2
(deficit 0.0400). The crossing is visible in the deficit tables without
additional diagnostics. [body, Section~\ref{sec:thirddoll}]

\medskip
\noindent\textbf{F7. Wiring succession occurs within the same model.}
The reverse-canonical wiring of cap4\_acc80 outperforms its canonical wiring
on the cap3 world (0.0577 vs 0.2405). On the cap4 world, the canonical wiring
wins (0.0400 vs 0.9823). The model remains; the wiring changes.
[body, Section~\ref{sec:thirddoll}]

\medskip
\noindent\textbf{F8. The correct successor was not the stage-1 leader.}
The acc80 canonical wiring that wins stage 2 (deficit 0.0400) had stage-1
deficit 0.3795 (rank 4). Committing to the stage-1 best would have yielded
deficit 10.31 at stage 2. [body, Section~\ref{sec:thirddoll}]

\subsection*{Signal structure findings}

\noindent\textbf{F9. World extension produces a directed,
boundary-concentrated signal.}
The ghost's deficit concentrates at C2 and C3 (98.9\% of total). The
contamination gradient originates at the untrained slot and decays with
distance from the boundary. Replicated at cap5 (C3/C4 boundary).
[body, Section~\ref{sec:twokinds}]

\medskip
\noindent\textbf{F10. Screen violation produces a uniform, inverted signal.}
Under 10\% count corruption, all wirings degrade. The ghost outperforms the
native (1.11 vs 1.69). Selecting on violation signal would promote the wrong
model. [body, Section~\ref{sec:twokinds}]

\medskip
\noindent\textbf{F11. The inversion holds across all contenders.}
Scoring all 18 stage-1 wirings under violation, ghost wirings dominate the top
ranks. The native acc80 canonical wiring --- the one that wins world extension
--- ranks 7th. The pattern replicates at cap5. [appendix]

\medskip
\noindent\textbf{F12. Lateral detection: model inadequacy is detectable before
new tokens arrive.}
Filtering C3 from the evaluation stream, the ghost still fails (deficit 3.3845
vs native's 0.0400). Boundary dynamics leak through INC/DEC before the
explicit new token appears. [body, Section~\ref{sec:twokinds}]

\subsection*{Implementation findings}

\noindent\textbf{F13. The interface principle is a functional requirement.}
Same model weights, two modes: bypass\_on scores 0.0518; bypass\_off scores
6.0289 with scrambled wiring. Violating the interface makes matching
unsolvable. [appendix]

\begin{center}
\begin{tabular}{lll}
\toprule
Mode & Best deficit & Best wiring \\
\midrule
bypass\_on  & \textbf{0.0518} & canonical \\
bypass\_off & 6.0289          & scrambled \\
\bottomrule
\end{tabular}
\end{center}

\medskip
\noindent\textbf{F14. Repeg degrades matching monotonically.}
Re-pegging the GRU hidden state toward trained count attractors during matching
hurts performance at every tested alpha value (0.0 $\to$ 0.1 $\to$ 0.2 $\to$
0.5: deficit 0.05 $\to$ 1.38 $\to$ 2.27 $\to$ 3.25). The model's internal
coordinate system is its own business. [appendix]

\medskip
\noindent\textbf{F15. Curriculum training is necessary for clean internal
geometry.}
Without staged curriculum, the GRU achieves similar accuracy ($\sim$80\%) but
with internally scrambled count state embeddings. The deficit landscape becomes
flatter and less discriminating. [appendix]

\medskip
\noindent\textbf{F16. All findings are robust across scoring functions.}
Ghost failure ordering and succession crossing hold under three scorers: deficit
(hinge), tiered (stepped penalty), and cross-entropy (NLL). The succession
event is in the structure, not in the scorer. [appendix]

\begin{center}
\begin{tabular}{llll}
\toprule
Model    & deficit & tiered & cross\_entropy \\
\midrule
symbolic & 0.0000  & 0.0000 & 283.7 \\
native   & 0.0400  & 1.80   & 382.1 \\
ghost    & 7.0694  & 219.2  & 537.9 \\
\bottomrule
\end{tabular}
\end{center}

\begin{center}
\begin{tabular}{lll}
\toprule
Scorer        & Stage-1 leader      & Stage-2 leader      \\
\midrule
deficit       & ghost (0.0000)      & native (0.0400)     \\
tiered        & ghost (0.0000)      & native (1.80)       \\
cross\_entropy & ghost (374.0)      & native (382.1)      \\
\bottomrule
\end{tabular}
\end{center}

Cross-entropy reveals what the deficit threshold hides: the symbolic model
(283.7) and native model (382.1) are distinguishable, and wrong wirings are
punished exponentially (871--998 vs 4.33--5.27 under deficit).

\medskip
\noindent\textbf{F17. Model quality affects discrimination sharpness.}
acc70 produces 8 stage-1 contenders vs acc80's 4. Lower quality yields a
flatter deficit landscape with smaller gaps. acc70's best stage-1 wiring is
shifted (non-canonical) --- a different training seed broke the symmetry
differently. [appendix]

\medskip
\noindent\textbf{F18. The contender pool is an existential requirement.}
Four pruning policies tested on cap5 three-stage succession: [appendix]

\begin{center}
\small
\begin{tabular}{lllll}
\toprule
Policy & Native contenders & Ghost contenders & Final winner & Final deficit \\
\midrule
threshold=0.05  & \textbf{0} & 1  & ghost            & 2.8887 \\
cliff (current) & 2          & 12 & \textbf{native}  & 0.0391 \\
threshold=2.0   & 4          & 4  & \textbf{native}  & 0.0391 \\
topk=1          & 1          & 0  & native$\dagger$  & 18.4069 \\
\bottomrule
\end{tabular}
\end{center}

Aggressive pruning (threshold=0.05) eliminates the correct model before the
decisive stage. Greedy commitment (topk=1) locks onto a path that yields
catastrophic deficit. The cliff policy succeeds because it reads the
distribution rather than imposing a fixed threshold. The contender set is the
architecture's hedge against world changes it has not yet observed.

\section{Experimental Conditions}
\label{app:conditions}

All experiments: n\_seeds=5, seq\_len=300, seed=42, p\_read=0.6. Deficit
scorer: $\max(1/n_{\mathrm{tokens}} - p_{\mathrm{pred}}(\mathrm{actual\_next}),
0)$ summed over steps, averaged over sequences.

\begin{center}
\resizebox{\textwidth}{!}{%
\begin{tabular}{llll}
\toprule
Experiment & Arena & Models & World \\
\midrule
Underdetermination   & cap3\_ghost\_games          & 3$\times$cap3, ghost, cap4\_acc80      & cap3 \\
GRU bypass           & cap4\_bypass\_compare       & cap4\_acc80 (two modes)                & cap4 \\
Repeg sweep          & cap4\_repeg\_sweep          & cap4\_acc80 (4 alphas), symbolic       & cap4 \\
Ghost (cap4)         & cap4\_ghost\_games          & ghost, native, symbolic                & cap4 \\
Ghost (cap5)         & cap5\_ghost\_games          & cap5\_ghost, cap5\_acc70, cap5\_acc80  & cap5 \\
Error structure      & cap4\_error\_structure      & ghost, native    & cap4 (A) / cap3+corrupt (B) \\
Error structure      & cap5\_error\_structure      & ghost, native    & cap5 (A) / cap4+corrupt (B) \\
Lateral detection    & cap4\_lateral\_detection    & ghost, native    & cap4 (filtered/unfiltered) \\
Succession (2-stage) & cap4\_succession            & ghost, cap4\_acc80   & cap3$\to$cap4 \\
Succession (3-stage) & cap5\_incremental           & ghost, cap5\_acc70, cap5\_acc80  & cap3$\to$cap4$\to$cap5 \\
Model quality        & cap4\_incremental\_topk     & ghost, acc70, acc80  & cap3$\to$cap4 \\
Scorer robustness    & cap4\_ghost\_games\_scorers & ghost, native, symbolic  & cap4 \\
Scorer robustness    & cap4\_succession\_scorers   & ghost, cap4\_acc80   & cap3$\to$cap4 \\
\bottomrule
\end{tabular}}
\end{center}

Full deficit tables for all experiments are available in the repository.

\section{Model Specifications}
\label{app:models}

All learned models share the same architecture: token embedding dimension 8,
GRU hidden dimension 16, one GRU layer. Token vocabulary: INC (0), DEC (1),
C0\ldots C$_{n-1}$ (2\ldots$n$+1), giving $n_{\mathrm{tokens}} =
n_{\mathrm{counts}} + 2$.

\begin{center}
\resizebox{\textwidth}{!}{%
\begin{tabular}{lrrrllr}
\toprule
Model & $n_{\mathrm{tok}}$ & $n_{\mathrm{cnt}}$ & Emb & GRU & Training & Accuracy \\
\midrule
cap3\_acc70              & 5 & 3 & 8 & 16 & Full cap3 curriculum     & $\sim$70\% \\
cap3\_acc80              & 5 & 3 & 8 & 16 & Full cap3 curriculum     & $\sim$80\% \\
cap4\_acc70              & 6 & 4 & 8 & 16 & Full cap4 curriculum     & $\sim$70\% \\
cap4\_acc80              & 6 & 4 & 8 & 16 & Full cap4 curriculum     & $\sim$80\% \\
cap4\_ghost              & 6 & 4 & 8 & 16 & Cap3 curriculum only     & $\sim$80\% on cap3 \\
cap5\_acc70              & 7 & 5 & 8 & 16 & Full cap5 curriculum     & $\sim$70\% \\
cap5\_acc80              & 7 & 5 & 8 & 16 & Full cap5 curriculum     & $\sim$80\% \\
cap5\_cap4trained\_acc80 & 7 & 5 & 8 & 16 & Cap4 curriculum only     & $\sim$80\% on cap4 \\
symbolic\_cap4           & 6 & 4 & --- & --- & By construction        & 100\% \\
\bottomrule
\end{tabular}}
\end{center}

\begin{thebibliography}{9}

\bibitem{wynn1992}
Wynn, K. (1992). Addition and subtraction by human infants. \textit{Nature},
358, 749--750.

\bibitem{architecture}
[Architecture paper --- arXiv ID to be assigned at simultaneous submission]

\bibitem{repo}
Repository: \url{https://github.com/kubasiemion/cognitive-architecture/tree/main/demos/wynns_world}

\end{thebibliography}

\end{document}
